%============================================================
\chapter{Lòng Tốt Và Nhân Từ (Kindness and Benevolence)}
\begin{center}
  \textit{\color{truyenblue}No act of kindness, no matter how small,\\
  is ever wasted.}\\
  \textit{(Không có hành động tốt bụng nào, dù nhỏ bé đến đâu,\\
  là bị lãng phí.)}\\[4pt]
  \small\color{truyengold}--- Aesop ---
\end{center}
\ngancach

\section{Mẹ Teresa Và Người Hấp Hối Bên Đường}
\begin{truyen}{Mẹ Teresa -- Chết Trong Nhân Phẩm}{Ấn Độ, Thế Kỷ XX}
\chuhoa{M}{}M ột ngày năm 1952, Mẹ Teresa nhặt được một người phụ nữ đang hấp hối bên cống rãnh Calcutta.\\
\textit{(One day in 1952, Mother Teresa found a woman dying in the gutters of Calcutta.)}

Chuột đang gặm chân bà, kiến đang bu quanh vết thương, không ai dừng lại nhìn.\\
\textit{(Rats were gnawing at her feet, ants were swarming her wounds, and no one stopped to look.)}

Mẹ Teresa bế bà vào bệnh viện, họ từ chối vì ``người sắp chết không đáng dùng giường bệnh.''\\
\textit{(Mother Teresa carried her to the hospital; they refused because ``the dying are not worth a hospital bed.'')}

Mẹ Teresa thuê một ngôi nhà, lập Nirmal Hriday -- ``Trái Tim Thuần Khiết'' -- nơi người hấp hối được chết trong nhân phẩm.\\
\textit{(Mother Teresa rented a building and established Nirmal Hriday -- ``Pure Heart'' -- where the dying could die with dignity.)}

Bà nói: ``Tôi không làm những việc lớn. Tôi chỉ làm những việc nhỏ với tình yêu lớn.''\\
\textit{(She said: ``I do not do great things. I only do small things with great love.'')}

\end{truyen}
\begin{baihoc}
  Lòng nhân từ không cần lý do hay phần thưởng. Mỗi người xứng đáng được đối xử với nhân phẩm -- dù ở đâu, lúc nào.\\
  \textit{(Kindness needs no reason or reward. Every person deserves to be treated with dignity -- wherever, whenever.)}
\end{baihoc}

\section{Truyện Người Samari Tốt Lành}
\begin{truyen}{Người Samari Tốt Lành}{Kinh Thánh -- Phúc Âm Luca 10:25--37}
\chuhoa{M}{}M ột người Do Thái bị cướp tấn công trên đường, bị đánh nằm chết dở bên vệ đường.\\
\textit{(A Jewish man was attacked by robbers on the road, beaten and left half-dead by the roadside.)}

Một thầy tế lễ đi qua, thấy rồi... tránh sang phía bên kia đường đi tiếp.\\
\textit{(A priest came by, saw him and... moved to the other side of the road and continued on.)}

Một người Lê-vi -- cũng người đạo đức -- cũng làm vậy.\\
\textit{(A Levite -- also a religious man -- did the same.)}

Rồi một người Samari -- bị người Do Thái coi thường -- dừng lại, băng bó vết thương, đưa vào quán trọ, trả tiền cho chủ quán chăm sóc.\\
\textit{(Then a Samaritan -- despised by Jews -- stopped, bandaged the wounds, brought him to an inn, and paid the innkeeper to care for him.)}

Jesus hỏi: ``Ai trong ba người đó đã là người láng giềng của người bị nạn?''\\
\textit{(Jesus asked: ``Which of these three do you think was a neighbor to the man who fell into the hands of robbers?'')}

\end{truyen}
\begin{baihoc}
  Lòng tốt không phân biệt chủng tộc, tôn giáo, địa vị. Người láng giềng là bất kỳ ai cần sự giúp đỡ của ta lúc này.\\
  \textit{(Kindness does not distinguish race, religion, or status. A neighbor is whoever needs our help at this moment.)}
\end{baihoc}

\section{Albert Schweitzer -- Bác Sĩ Bỏ Danh Tiếng Vào Rừng Châu Phi}
\begin{truyen}{Albert Schweitzer -- Sứ Mệnh Tình Người}{Châu Phi, Thế Kỷ XX}
\chuhoa{A}{}A lbert Schweitzer ở tuổi 30 là nhà triết học, nhạc sĩ organ hàng đầu châu Âu, giáo sư thần học nổi tiếng.\\
\textit{(At age 30, Albert Schweitzer was a leading philosopher, organist, and renowned theology professor in Europe.)}

Ông đọc bài báo về dịch bệnh tàn phá vùng Gabon, châu Phi -- nơi không có bác sĩ.\\
\textit{(He read an article about a devastating epidemic in Gabon, Africa -- where there were no doctors.)}

Ông quyết định học y khoa từ đầu và lên đường sang châu Phi xây bệnh viện trong rừng.\\
\textit{(He decided to study medicine from scratch and set off to Africa to build a hospital in the jungle.)}

Từ bỏ danh tiếng và tiện nghi, ông sống và làm việc ở Lambaréné hơn 50 năm.\\
\textit{(Abandoning his fame and comforts, he lived and worked in Lambarene for over 50 years.)}

Khi nhận Nobel Hòa Bình năm 1952, ông dùng toàn bộ tiền thưởng xây thêm bệnh viện phong.\\
\textit{(When he received the Nobel Peace Prize in 1952, he used the entire prize money to build a leprosy hospital.)}

\end{truyen}
\begin{baihoc}
  Lòng nhân từ cao nhất là trao đi những điều quý giá nhất của mình -- thời gian, tài năng, cơ hội -- cho người cần nhất.\\
  \textit{(The highest form of kindness is giving the most precious things you have -- time, talent, opportunity -- to those who need it most.)}
\end{baihoc}

\section{Câu Chuyện Những Tờ Giấy Khen Ngợi}
\begin{truyen}{Thí Nghiệm Về Lòng Tốt}{Câu Chuyện Giáo Dục Hiện Đại}
\chuhoa{M}{}M ột giáo viên tiểu học yêu cầu học sinh viết một điều tốt đẹp nhất về mỗi bạn trong lớp.\\
\textit{(An elementary school teacher asked students to write one best thing about each classmate.)}

Cô thu lại, tổng hợp và trao cho từng em tờ giấy ghi những điều bạn bè nói về mình.\\
\textit{(She collected them, compiled the comments, and gave each child a paper listing what their friends said about them.)}

Mắt các em sáng lên -- nhiều em lần đầu biết rằng mình được bạn bè yêu quý như vậy.\\
\textit{(The children's eyes lit up -- many learned for the first time how much their friends valued them.)}

Nhiều năm sau, một học sinh cũ -- lính chiến đã mất -- được cha mẹ mang về tờ giấy ố vàng đó trong túi áo.\\
\textit{(Years later, a former student -- a fallen soldier -- was brought back by his parents with that yellowed paper still in his breast pocket.)}

Nhiều học trò khác cũng thú nhận: họ vẫn còn giữ tờ giấy đó sau bao nhiêu năm.\\
\textit{(Other former students also confessed: they had kept their papers all those years.)}

\end{truyen}
\begin{baihoc}
  Lời khen ngợi chân thành dù nhỏ bé có thể trở thành ánh sáng soi đường cho người khác trong những tháng năm tối tăm nhất.\\
  \textit{(A sincere word of praise, however small, can become a guiding light for others in their darkest times.)}
\end{baihoc}

\section{Phạm Ngũ Lão Đan Kiếm Chờ Thời}
\begin{truyen}{Phạm Ngũ Lão Ngồi Đan Kiếm Giữa Đường}{Việt Nam, Thế Kỷ XIII}
\chuhoa{P}{}P hạm Ngũ Lão -- chàng trai nghèo -- đang ngồi đan chiếu cói giữa đường làng.\\
\textit{(Pham Ngu Lao, a poor young man, was sitting weaving a rush mat in the middle of the village road.)}

Đoàn quân của Hưng Đạo Vương đi qua, quân lính hét to nhưng anh ngồi yên không nhúc nhích.\\
\textit{(The troops of Prince Hung Dao passed by; soldiers shouted but he sat still without moving.)}

Một ngọn giáo đâm vào đùi anh -- anh vẫn không rời tầm mắt khỏi chiếc chiếu đang đan dở.\\
\textit{(A spear was thrust into his thigh -- he still did not take his eyes off the unfinished mat.)}

Hưng Đạo Vương lạ lùng dừng lại hỏi: anh đang suy nghĩ điều gì mà phép thần như vậy?\\
\textit{(Prince Hung Dao, astonished, stopped and asked: what were you thinking to be so tranquil?)}

Phạm Ngũ Lão thưa: ``Thần đang suy nghĩ kế sách phá giặc Nguyên cho đất nước.''\\
\textit{(Pham Ngu Lao replied: ``I was devising a strategy to defeat the Yuan invaders for our country.'')}

Hưng Đạo Vương lập tức thu nhận và Phạm Ngũ Lão trở thành danh tướng bách chiến bách thắng.\\
\textit{(Prince Hung Dao immediately recruited him and Pham Ngu Lao became a general who won every battle.)}

\end{truyen}
\begin{baihoc}
  Lòng tốt với đất nước và nhân dân -- thể hiện qua sự tận tâm âm thầm -- bao giờ cũng được nhận ra bởi những người xứng đáng.\\
  \textit{(Devotion to one's country and people -- expressed through quiet dedication -- is always recognized by those worthy of it.)}
\end{baihoc}
